\documentclass[11pt]{article}

\usepackage[letterpaper,margin=1in]{geometry}
\usepackage{amsmath}
\usepackage{amssymb}
\usepackage{array}
\usepackage{booktabs}
\usepackage{longtable}
\usepackage{hyperref}
\usepackage{xcolor}
\usepackage{listings}
\usepackage{fancyhdr}

\hypersetup{
  colorlinks=true,
  linkcolor=blue,
  citecolor=blue,
  urlcolor=blue,
  pdftitle={PADZE User Manual},
  pdfauthor={Andres del Castillo and Yitzchak Shmalo}
}

% ---- Listing styles -------------------------------------------------------
\lstset{
  breaklines=true,
  breakatwhitespace=false,
  columns=fullflexible,
  frame=single,
  framerule=0.2pt,
  xleftmargin=0.6em,
  xrightmargin=0.6em,
  aboveskip=0.9em,
  belowskip=0.9em,
  showstringspaces=false,
  keepspaces=true
}
\lstdefinestyle{cmd}{basicstyle=\ttfamily\small,language=bash}
\lstdefinestyle{py}{basicstyle=\ttfamily\small,language=Python}
\lstdefinestyle{out}{basicstyle=\ttfamily\footnotesize}

% ---- "Coming from ADZE" callout box ---------------------------------------
\definecolor{calloutbg}{RGB}{240,244,248}
\definecolor{calloutrule}{RGB}{71,85,105}
\newsavebox{\calloutbox}
\newenvironment{callout}[1]{%
  \par\medskip\noindent
  \setlength{\fboxsep}{10pt}%
  \setlength{\fboxrule}{0.8pt}%
  \begin{lrbox}{\calloutbox}%
  \begin{minipage}{\dimexpr\textwidth-2\fboxsep-2\fboxrule\relax}%
    \small
    \textbf{\large #1}\par\smallskip
}{%
  \end{minipage}%
  \end{lrbox}%
  \noindent\fcolorbox{calloutrule}{calloutbg}{\usebox{\calloutbox}}%
  \par\medskip
}

\pagestyle{fancy}
\setlength{\headheight}{14pt}
\fancyhf{}
\lhead{PADZE User Manual}
\rhead{Version 0.1.0}
\cfoot{\thepage}

\newcommand{\padze}{PADZE}
\setcounter{tocdepth}{2}
\setcounter{secnumdepth}{2}

% Let TeX absorb the last few points of long inline flag names in justified
% prose instead of leaving them slightly past the right margin.
\setlength{\emergencystretch}{3em}

\begin{document}
\pagenumbering{roman}

% ===========================================================================
\begin{titlepage}
\centering
\vspace*{1.25in}
{\Huge \textbf{\padze{}}\par}
\vspace{0.25in}
{\Large Pythonic Allelic Diversity Analyzer\par}
\vspace{0.25in}
{\large User Manual \textbullet{} Version 0.1.0\par}
\vspace{0.75in}
{\large Andres del Castillo and Yitzchak Shmalo\par}
\vspace{0.35in}
{\large \today\par}
\vspace{0.15in}
{\large Distributed under the MIT license\par}
\vfill
\begin{minipage}{0.86\textwidth}
\padze{} (Pythonic Allelic Diversity Analyzer) measures how many distinct
alleles each of several populations carries, and corrects the counts by
rarefaction so that populations sampled at different sizes can be compared
fairly. For each locus and standardized sample depth it reports allelic
richness, private allelic richness, and richness private to a named
combination of populations, and then summarizes each statistic across loci.
\padze{} reads VCF or STRUCTURE genotypes, writes tidy CSV or ready-to-use
summary files, and runs from either a \texttt{padze} command line or an
importable Python API.
\end{minipage}
\vfill
\begin{minipage}{0.86\textwidth}
\small
\padze{} implements the rarefaction estimators of Szpiech, Jakobsson \&
Rosenberg (2008); cite that paper when you report results. The full citation is
in the References.
\end{minipage}
\end{titlepage}

\setcounter{page}{2}
\tableofcontents
\newpage
\pagenumbering{arabic}

% ===========================================================================
\section{Introduction}

The number of distinct alleles a population carries---its \emph{allelic
diversity}---is a basic measure of genetic variation. It is also treacherous to
compare across populations, because a larger sample turns up more alleles simply
by examining more gene copies. \padze{} removes that sampling bias by
\emph{rarefaction}: for every population it estimates how many alleles you would
expect to see if you had drawn the same standardized number of gene copies from
each. The corrected counts are directly comparable even when the raw sample
sizes are not.

For each locus and each standardized sample depth \(g\), \padze{} reports three
quantities:

\begin{itemize}
  \item \textbf{Allelic richness} (\texttt{alpha})---expected distinct alleles
        in a population.
  \item \textbf{Private allelic richness} (\texttt{pi})---expected alleles
        unique to one population.
  \item \textbf{Combination-private richness} (\texttt{pihat})---expected
        alleles unique to a named \emph{group} of populations.
\end{itemize}

It computes these at every locus and then summarizes each statistic
\emph{across loci}, reporting the mean, variance, and standard error together
with the skewness and excess kurtosis of the across-loci distribution. The
result is a compact numerical profile of how diversity is distributed across the
genome, in a CSV you can analyze directly or feed to downstream models.

\padze{} is packaged as the Python project \texttt{padze}, with console command
\texttt{padze} and module entry point \texttt{python -m padze}. It is
distributed under the MIT license.

\begin{callout}{Coming from ADZE 1.0?}
\padze{} runs your existing analyses and adds new capability:
\begin{itemize}
  \setlength{\itemsep}{2pt}
  \item \textbf{Input.} Point \padze{} at a VCF plus a one-line-per-sample
        population map; no STRUCTURE conversion. Your existing STRUCTURE files
        still work.
  \item \textbf{Install and run.} Install with \texttt{pip} on any platform,
        then use the \texttt{padze} CLI or the Python API, in place of
        per-platform binaries.
  \item \textbf{Reproduce R/P/C.} Run your data with \texttt{-{}-adze-prefix} to
        regenerate the R, P, and C files (and the FULL per-locus files)
        directly.
  \item \textbf{New statistics.} The same mean, variance, and standard error,
        plus across-loci skewness and excess kurtosis, and classical statistics
        over sliding genomic windows.
\end{itemize}
\padze{} matches ADZE 1.0 wherever they overlap. Full mapping and parity
details are in \texttt{docs/adze-1.0-reference.md}.
\end{callout}

% ===========================================================================
\section{Core Concepts}
\label{sec:concepts}

This section explains everything needed to read \padze{}'s output. If you
already work with rarefied allelic diversity, skim to
Section~\ref{sec:depth} for the depth rule and Section~\ref{sec:defs} for the
formulas.

\subsection{Rarefaction and Gene Copies}

Diversity statistics are counted in \emph{gene copies}, not individuals. A
diploid individual contributes two gene copies per locus, a haploid one, and so
on. Two diploid samples therefore provide up to four gene copies at a locus.

Because a bigger sample finds more alleles just by looking harder, raw counts are
not comparable across populations of different size. \emph{Rarefaction} fixes
this by trimming every population to a common standardized depth \(g\) (a number
of gene copies) and asking: \emph{if you had drawn only \(g\) gene copies from
this population, how many distinct alleles would you expect to see?} Answering
that question for a range of depths \(g = 2, 3, \dots\) is the core of what
\padze{} does. Combinatorial formulas give the exact expectation, so no
resampling is needed.

\subsection{The Three Richness Statistics}

At a single locus and depth \(g\), \padze{} evaluates:

\begin{itemize}
  \item \texttt{alpha\_j}---\textbf{allelic richness.} Expected number of
        distinct alleles observed when \(g\) gene copies are drawn from
        population \(j\).
  \item \texttt{pi\_j}---\textbf{private allelic richness.} Expected number of
        alleles \emph{private} to population \(j\)---present in \(j\) and absent
        from every other population---at depth \(g\).
  \item \texttt{pihat}---\textbf{combination-private richness.} Expected number
        of alleles private to a \emph{combination} of populations: present in
        every member of the named group and absent everywhere outside it.
        \texttt{pihat} generalizes \texttt{pi} (a single-population combination
        is exactly \texttt{pi}), which is why it carries the ``hat.'' For
        example \texttt{pihat\_12} counts alleles found in populations 1 and 2
        but in no other.
\end{itemize}

\emph{Private} always means present here, absent everywhere else. The numeric
suffix indexes populations in the population order (by default, the order in
which they first appear), so with \(P\) populations in total \texttt{alpha\_1}
is the first population, \texttt{alpha\_2} the second, and \texttt{pihat\_13} is
the combination of the first and third. In the shipped example (populations
\(A, B, C\)) that makes \texttt{alpha\_1} population \(A\) and \texttt{pihat\_13}
the \(A\)--\(C\) combination.

\subsection{Summaries Across Loci}

Each statistic has one value \emph{per locus}. Across the loci in a dataset those
values form a distribution, and \padze{} summarizes its shape at each depth with
five moments:

\begin{center}
\begin{tabular}{@{}>{\raggedright\arraybackslash}p{0.20\textwidth}>{\raggedright\arraybackslash}p{0.72\textwidth}@{}}
\toprule
\textbf{Moment} & \textbf{Meaning} \\
\midrule
\texttt{mean} & Average across loci---the headline diversity estimate. \\
\texttt{variance} & Sample variance across loci (locus-to-locus spread). \\
\texttt{se} & Standard error of the mean, \(\sqrt{\text{variance}/n_{\text{loci}}}\). \\
\texttt{skewness} & Asymmetry of the across-loci distribution. \\
\texttt{kurtosis} & Excess kurtosis---tailedness relative to a normal distribution. \\
\bottomrule
\end{tabular}
\end{center}

The \texttt{mean}, \texttt{variance}, and \texttt{se} are the classical
summaries. The \texttt{skewness} and \texttt{kurtosis} add a description of the
distribution's asymmetry and tails---useful when diversity is concentrated in a
few loci rather than spread evenly.

\subsection{Rarefaction Depth}
\label{sec:depth}

You can only rarefy down to a depth the data actually support. The maximum usable
depth is set by the \textbf{smallest number of non-missing gene copies}
available for a population at any single locus. Two diploid samples give four
gene copies per population, but a single missing genotype leaves three copies at
that locus---so the depth for that population is capped at three.

Choosing \texttt{-{}-max-depth} is a trade-off: a larger \(g\) sits closer to the
raw allele count but is supported by fewer loci and populations, while a smaller
\(g\) is more comparable across small samples. \padze{} reports the maximum depth
your data support in the \texttt{info} summary (Section~\ref{sec:info}).

\subsection{Precise Definitions}
\label{sec:defs}

Consider a locus with \(I\) distinct alleles. Let \(N_{ij}\) be the number of
copies of allele \(i\) in population \(j\), and
\[
  N_j = \sum_{i=1}^{I} N_{ij}
\]
the total number of gene copies in \(j\). The probability that a draw of \(g\)
gene copies from \(j\) (without replacement) misses allele \(i\) entirely is
\[
  Q_{ijg} = \frac{\dbinom{N_j - N_{ij}}{g}}{\dbinom{N_j}{g}},
\]
and \(P_{ijg} = 1 - Q_{ijg}\) is the probability the draw contains allele \(i\).
Summing over the alleles at the locus,
\[
  \hat{\alpha}^{(j)}_g = \sum_{i=1}^{I} P_{ijg},
\qquad
  \hat{\pi}^{(j)}_g = \sum_{i=1}^{I} P_{ijg}\!\!\prod_{\substack{j'=1\\ j' \ne j}}^{P}\!\! Q_{ij'g},
\]
\[
  \hat{\pi}^{(S)}_g = \sum_{i=1}^{I}
  \Bigl(\prod_{j \in S} P_{ijg}\Bigr)
  \Bigl(\prod_{j \notin S} Q_{ijg}\Bigr).
\]
Here \(\hat{\alpha}^{(j)}_g\) counts alleles the draw is expected to contain;
\(\hat{\pi}^{(j)}_g\) counts alleles present in \(j\) but absent from every other
population; and \(\hat{\pi}^{(S)}_g\) counts alleles present in every population
of the set \(S\) and absent from all populations outside it. Across \(L\) loci,
the summaries of a per-locus value \(x\) are
\[
  \bar{x} = \frac{1}{L}\sum_{\ell=1}^{L} x_{\ell},
\qquad
  s^2 = \frac{1}{L-1}\sum_{\ell=1}^{L}(x_{\ell}-\bar{x})^2,
\qquad
  \mathrm{SE} = \sqrt{\frac{s^2}{L}}.
\]

% ===========================================================================
\section{Installation}
\label{sec:install}

\padze{} requires Python 3.10 or newer and NumPy. The optional \texttt{test}
extra adds pytest and SciPy for the validation suite.

Install from a checkout of the repository:
\begin{lstlisting}[style=cmd]
pip install .
\end{lstlisting}

Install with the test/validation dependencies:
\begin{lstlisting}[style=cmd]
pip install -e ".[test]"
\end{lstlisting}

Installing creates the \texttt{padze} console command; you can also invoke the
package as a module:
\begin{lstlisting}[style=cmd]
padze --help
python -m padze --help
\end{lstlisting}

To run \emph{without installing}, from the \texttt{GitHub/} directory of the
checkout put the source on the path and call the module:
\begin{lstlisting}[style=cmd]
PYTHONPATH=src python -m padze --help
\end{lstlisting}

Every command in this manual uses the no-install form
\texttt{PYTHONPATH=src python -m padze \ldots} so it runs directly from a fresh
checkout. After \texttt{pip install .} you can drop the
\texttt{PYTHONPATH=src python -m} prefix and just type \texttt{padze}.

% ===========================================================================
\section{Quick Start}
\label{sec:quickstart}

\padze{} ships a small example in \texttt{examples/}: \texttt{trio.vcf} and
\texttt{trio.popmap} describe three populations (\(A, B, C\)), two diploid
samples each, across seven loci.

\subsection{Inspect the Data with \texttt{info}}
\label{sec:info}

Start with \texttt{info} to confirm \padze{} reads your data as you expect:
\begin{lstlisting}[style=cmd]
PYTHONPATH=src python -m padze info \
  --vcf examples/trio.vcf \
  --popmap examples/trio.popmap
\end{lstlisting}

Output:
\begin{lstlisting}[style=out]
source: VCF:examples/trio.vcf
populations (3): A, B, C
samples/pop: A=2, B=2, C=2
loci: kept 7 / read 7
filters: drop all-missing loci
missing call fraction (kept): 0.0119
ADZE feature shape per locus: alpha_j, pi_j (j=1..P) and pihat over population
combinations; summarized across loci as (mean, variance, se, skewness,
kurtosis) per rarefaction depth.
max usable rarefaction depth: 3
\end{lstlisting}

The depth caps at 3 (not 4) because a missing genotype at \texttt{chr1:600}
leaves population \(A\) with three gene copies at that locus---see
Section~\ref{sec:depth}.

\subsection{Compute the Feature Table}
\label{sec:qs-features}

The default \texttt{features} output is a wide CSV: one row per rarefaction depth
\(g\), with the five moment columns for every statistic. Use \texttt{-{}-out -}
to write to standard output.
\begin{lstlisting}[style=cmd]
PYTHONPATH=src python -m padze features \
  --vcf examples/trio.vcf \
  --popmap examples/trio.popmap \
  --max-depth 3 \
  --out -
\end{lstlisting}

Output (columns and rows truncated; 46 columns total, full schema in
Section~\ref{sec:outputs}):
\begin{lstlisting}[style=out]
g,alpha_1_mean,alpha_1_variance,alpha_1_se,alpha_1_skewness,alpha_1_kurtosis, ... ,pihat_23_kurtosis
2,1.380952380952381,0.07142857142857144,0.10101525445522108,-0.9839173128183324, ...
3,1.5714285714285714,0.16071428571428573,0.15152288168283162,-0.983917312818333, ...
\end{lstlisting}

Read the first row as: population \(A\) (\texttt{alpha\_1}), rarefied to
\(g = 2\) gene copies and averaged over 7 loci, is expected to carry
\textbf{1.38 distinct alleles} per locus, with variance 0.071 and standard error
0.101 across those loci.

\subsection{Compute the Classical Summary}
\label{sec:qs-classical}

Pass \texttt{-{}-classical} for a compact long-format table of just the three
classical moments---mean, variance, and standard error---one row per statistic
and depth:
\begin{lstlisting}[style=cmd]
PYTHONPATH=src python -m padze features \
  --vcf examples/trio.vcf \
  --popmap examples/trio.popmap \
  --max-depth 3 \
  --classical \
  --out -
\end{lstlisting}

Output:
\begin{lstlisting}[style=out]
statistic,g,n_loci,mean,variance,se
alpha_1,2,7,1.380952380952381,0.07142857142857144,0.10101525445522108
alpha_1,3,7,1.5714285714285714,0.16071428571428573,0.15152288168283162
alpha_2,2,7,1.3095238095238095,0.08730158730158728,0.11167656571008164
alpha_2,3,7,1.4642857142857144,0.19642857142857148,0.1675148485651225
alpha_3,2,7,1.476190476190476,0.05026455026455028,0.0847387162859628
alpha_3,3,7,1.7142857142857142,0.11309523809523812,0.1271080744289442
pi_1,2,7,0.047619047619047616,0.004298941798941799,0.02478173808888253
pi_1,3,7,0,0,0
pi_2,2,7,0.1547619047619048,0.031084656084656086,0.06663831596724867
pi_2,3,7,0.044642857142857144,0.008742559523809524,0.035340303830469995
pi_3,2,7,0.16666666666666669,0.13888888888888892,0.14085904245475278
pi_3,3,7,0.14285714285714285,0.14285714285714285,0.14285714285714285
pihat_12,2,7,0.32142857142857145,0.026455026455026454,0.061475926130276456
pihat_12,3,7,0.24107142857142855,0.06305803571428571,0.0949120161851308
pihat_13,2,7,0.4761904761904762,0.13128306878306875,0.13694788830743965
pihat_13,3,7,0.39285714285714285,0.18452380952380953,0.16235930591649828
pihat_23,2,7,0.2976190476190476,0.0853174603174603,0.11040022018447265
pihat_23,3,7,0.24107142857142858,0.12555803571428573,0.13392857142857142
\end{lstlisting}

\subsection{Python API}
\label{sec:api}

The same pipeline is available as an importable library:
\begin{lstlisting}[style=py]
import padze

loci = padze.read_vcf("examples/trio.vcf", "examples/trio.popmap")
table = padze.compute_features(loci)
matrix, columns = table.to_frame()
print(matrix.shape)        # (2, 46)
print(columns[:4])         # ['g', 'alpha_1_mean', 'alpha_1_variance', 'alpha_1_se']
\end{lstlisting}

% ===========================================================================
\section{Input Data}
\label{sec:input}

\padze{} accepts two input formats: \textbf{VCF plus a population map}
(recommended for new work) and \textbf{STRUCTURE} (for parity with, and
migration from, existing pipelines).

\subsection{VCF and Population Map}
\label{sec:vcf}

VCF is the primary input. \padze{} reads genotypes from the \texttt{GT} FORMAT
subfield:
\begin{itemize}
  \item Genotypes are taken from the \texttt{GT} subfield (the first
        \texttt{FORMAT} field, as usual).
  \item Ploidy is detected per genotype from the \texttt{GT} separators; samples
        are diploid by default, and a \texttt{.} allele is treated as a missing
        gene copy.
  \item Multiallelic sites are supported: alleles are aligned by VCF allele index
        across populations, so no biallelic decomposition is required.
  \item Plain \texttt{.vcf} and gzip-compressed \texttt{.vcf.gz} files are both
        accepted.
\end{itemize}

The \textbf{population map} (popmap) assigns each VCF sample to a population. It
is a plain-text file with one \texttt{sample<TAB>population} pair per line (any
whitespace is accepted; a tab is recommended). Lines beginning with \texttt{\#}
are comments. The shipped \texttt{examples/trio.popmap}:
\begin{lstlisting}[style=out]
# sample population
sA1	A
sA2	A
sB1	B
sB2	B
sC1	C
sC2	C
\end{lstlisting}

Population labels are ordered as first encountered. Use
\texttt{-{}-population-order} when the numbered outputs (\texttt{alpha\_1},
\texttt{pihat\_13}, \ldots) must follow a specific order, for example
\texttt{-{}-population-order A B C}.

\subsection{STRUCTURE Input}
\label{sec:structure}

STRUCTURE-format input lets you run the exact genotype files used by earlier
pipelines through \padze{}. A STRUCTURE file has an optional header row of locus
names, one or more leading metadata columns (individual label, population,
\ldots), and then the allele tokens. \padze{} reads both the
two-rows-per-individual layout (one row per gene copy) and the packed
one-row-per-individual layout. The relevant flags are documented in
Section~\ref{sec:flags}; a typical diploid, packed file with one header row and
two metadata columns (label, population) is read with:
\begin{lstlisting}[style=cmd]
PYTHONPATH=src python -m padze features \
  --structure mydata.stru \
  --structure-header-rows 1 \
  --structure-meta-columns 2 \
  --structure-pop-column 2 \
  --ploidy 2 \
  --one-row-per-individual \
  --population-order A B C \
  --max-depth 3 \
  --classical --out -
\end{lstlisting}

A STRUCTURE encoding of the trio example produces byte-for-byte the same
allelic-richness, private-richness, and \texttt{pihat} values as the matched VCF
(Section~\ref{sec:validation}). For the crosswalk between these flags and the
original ADZE parameter-file keys, see \texttt{docs/adze-1.0-reference.md}.

% ===========================================================================
\section{Running \padze{}}
\label{sec:running}

\subsection{Subcommands}
\label{sec:subcommands}

\padze{} has two subcommands, both taking the same input and filtering flags:
\begin{itemize}
  \item \texttt{info} loads the input, applies any filters, and prints a
        transparent summary---source, population labels, samples per population,
        loci read vs.\ kept, filters applied, missing-call fraction, and the
        maximum usable rarefaction depth.
  \item \texttt{features} runs the full rarefaction pipeline and writes the
        feature table (default), or the classical, ADZE-compatible, or
        rolling-window outputs on request.
\end{itemize}
\begin{lstlisting}[style=cmd]
PYTHONPATH=src python -m padze info     --vcf examples/trio.vcf --popmap examples/trio.popmap
PYTHONPATH=src python -m padze features --vcf examples/trio.vcf --popmap examples/trio.popmap
\end{lstlisting}

\subsection{Flag Reference}
\label{sec:flags}

\subsubsection*{Input and formats}
\begin{longtable}{@{}>{\raggedright\arraybackslash}p{0.38\textwidth}>{\raggedright\arraybackslash}p{0.58\textwidth}@{}}
\toprule
\textbf{Flag} & \textbf{Meaning} \\
\midrule
\endfirsthead
\toprule
\textbf{Flag} & \textbf{Meaning} \\
\midrule
\endhead
\texttt{-{}-vcf VCF} & Input \texttt{.vcf} / \texttt{.vcf.gz} file. \\
\texttt{-{}-popmap POPMAP} & Population map: one \texttt{sample population} pair per line; \texttt{\#} comments allowed. \\
\texttt{-{}-population-order P1 P2 \ldots} & Population order for the numbered outputs (\texttt{alpha\_1}, \texttt{pihat\_12}, \ldots). Default: order of first appearance. \\
\texttt{-{}-structure STRUCTURE} & Input STRUCTURE-format file. \\
\texttt{-{}-structure-header-rows N} & Leading non-data rows to skip (default 0). \\
\texttt{-{}-structure-meta-columns N} & Leading metadata columns before the genotypes (default: enough to include the label and population columns). \\
\texttt{-{}-structure-pop-column N} & 1-based column that assigns each individual to a population (default 2). \\
\texttt{-{}-structure-label-column N} & 1-based individual-label column (default 1). \\
\texttt{-{}-structure-missing TOKEN} & STRUCTURE missing-data token (default \texttt{-9}). \\
\texttt{-{}-ploidy P} & STRUCTURE ploidy (default 2). \\
\texttt{-{}-one-row-per-individual} & STRUCTURE rows pack \texttt{ploidy} alleles per locus (one row per individual). \\
\texttt{-{}-default-ploidy P} & VCF fallback ploidy for a bare missing \texttt{GT} \texttt{.} before a sample's ploidy is observed (default 2). \\
\bottomrule
\end{longtable}

\subsubsection*{Filtering and missingness}
\begin{longtable}{@{}>{\raggedright\arraybackslash}p{0.38\textwidth}>{\raggedright\arraybackslash}p{0.58\textwidth}@{}}
\toprule
\textbf{Flag} & \textbf{Meaning} \\
\midrule
\endfirsthead
\toprule
\textbf{Flag} & \textbf{Meaning} \\
\midrule
\endhead
\texttt{-{}-require-pass} & VCF: keep only records whose \texttt{FILTER} is \texttt{PASS} or \texttt{.}. \\
\texttt{-{}-max-missing-fraction F} & VCF: drop sites whose overall missing-call fraction exceeds \texttt{F}. \\
\texttt{-{}-max-missing-per-population F} & Drop a locus if \emph{any} population exceeds missing fraction \texttt{F}. \\
\texttt{-{}-min-populations-genotyped N} & Drop sites genotyped in fewer than \texttt{N} populations (default: all populations for the feature table, 1 for classical / ADZE-compatible modes). \\
\texttt{-{}-biallelic-only} & VCF: drop multiallelic sites. \\
\bottomrule
\end{longtable}

\subsubsection*{Statistics and depth}
\begin{longtable}{@{}>{\raggedright\arraybackslash}p{0.38\textwidth}>{\raggedright\arraybackslash}p{0.58\textwidth}@{}}
\toprule
\textbf{Flag} & \textbf{Meaning} \\
\midrule
\endfirsthead
\toprule
\textbf{Flag} & \textbf{Meaning} \\
\midrule
\endhead
\texttt{-{}-max-depth G} & Largest rarefaction depth \(g\) (default: the maximum the data support). \\
\texttt{-{}-pihat-sizes K [K \ldots]} & Combination sizes used for \texttt{pihat} (default: 2 when \(P \ge 2\), 1 when \(P = 1\)). \\
\texttt{-{}-moments M \ldots} & Which across-loci moments to write in the feature table; each \texttt{M} is one of \texttt{mean}, \texttt{variance}, \texttt{se}, \texttt{skewness}, \texttt{kurtosis} (default: all five). \\
\texttt{-{}-depth-policy \{common,ragged\}} & Feature-table depth handling: \texttt{common} requires every locus/population to support each depth; \texttt{ragged} summarizes whatever per-locus values are available up to the largest sample size. \\
\texttt{-{}-population-moments} & Use biased population skew/kurtosis (\texttt{g1}, \texttt{g2}) instead of the sample estimators (\texttt{G1}, \texttt{G2}). \\
\bottomrule
\end{longtable}

\subsubsection*{Output selection}
\begin{longtable}{@{}>{\raggedright\arraybackslash}p{0.38\textwidth}>{\raggedright\arraybackslash}p{0.58\textwidth}@{}}
\toprule
\textbf{Flag} & \textbf{Meaning} \\
\midrule
\endfirsthead
\toprule
\textbf{Flag} & \textbf{Meaning} \\
\midrule
\endhead
\texttt{-{}-out PATH} & Output CSV path (\texttt{-} writes to standard output). \\
\texttt{-{}-classical} & Emit the classical \texttt{(mean, variance, se)} table over each statistic's full depth range. \\
\texttt{-{}-adze-prefix PREFIX} & Write ADZE-compatible R/P/C files using this prefix (\texttt{PREFIX\_r}, \texttt{PREFIX\_p}, \texttt{PREFIX\_c\_K}). \\
\texttt{-{}-adze-full} & With \texttt{-{}-adze-prefix}, also write the FULL per-locus files. \\
\texttt{-{}-per-locus-out PATH} & Also write per-locus values in long format. \\
\bottomrule
\end{longtable}

\subsubsection*{Rolling windows}

Rolling windows compute the classical statistics over sliding stretches of the
genome, so you can trace how diversity changes along a chromosome.
\begin{longtable}{@{}>{\raggedright\arraybackslash}p{0.38\textwidth}>{\raggedright\arraybackslash}p{0.58\textwidth}@{}}
\toprule
\textbf{Flag} & \textbf{Meaning} \\
\midrule
\endfirsthead
\toprule
\textbf{Flag} & \textbf{Meaning} \\
\midrule
\endhead
\texttt{-{}-rolling-window W} & Compute the classical statistics over sliding windows of size \texttt{W}. \\
\texttt{-{}-step S} & Rolling-window step (default: \texttt{W}, i.e.\ non-overlapping windows). \\
\texttt{-{}-window-unit \{loci,bp\}} & Rolling-window unit: number of loci (default) or base pairs. \\
\bottomrule
\end{longtable}

Example---non-overlapping-ish windows of four loci, stepping by three:
\begin{lstlisting}[style=cmd]
PYTHONPATH=src python -m padze features \
  --vcf examples/trio.vcf \
  --popmap examples/trio.popmap \
  --max-depth 2 \
  --classical \
  --rolling-window 4 --step 3 --window-unit loci \
  --out -
\end{lstlisting}

Output (first rows):
\begin{lstlisting}[style=out]
window,unit,start,end,n_loci,statistic,g,n_used,mean,variance,se
0,loci,0,4,4,alpha_1,2,4,1.5,0,0
0,loci,0,4,4,alpha_2,2,4,1.25,0.08333333333333333,0.14433756729740643
0,loci,0,4,4,alpha_3,2,4,1.4583333333333335,0.09953703703703705,0.15774745405000762
0,loci,0,4,4,pi_1,2,4,0.0625,0.006365740740740739,0.03989279615651408
0,loci,0,4,4,pi_2,2,4,0.22916666666666669,0.03877314814814815,0.098454492213596
\end{lstlisting}

\subsection{Performance}
\label{sec:performance}

The rarefaction core is NumPy-vectorized and processes hundreds of loci and
thousands of individuals in seconds. The one cost to watch is combination-private
richness: the number of population combinations of size \(k\) grows as
\(\binom{P}{k}\) in the number of populations \(P\), so a large combination size
over many populations (set with \texttt{-{}-pihat-sizes}) can make the
\texttt{pihat} calculation dominate the run.

% ===========================================================================
\section{Output Files}
\label{sec:outputs}

\padze{} writes one of several formats depending on the flags you pass. Choose by
goal:

\begin{center}
\begin{tabular}{@{}>{\raggedright\arraybackslash}p{0.40\textwidth}>{\raggedright\arraybackslash}p{0.28\textwidth}>{\raggedright\arraybackslash}p{0.26\textwidth}@{}}
\toprule
\textbf{Goal} & \textbf{Flags} & \textbf{Format} \\
\midrule
Analyze in Python / feed downstream models & \emph{(default)} & Wide feature CSV (\S\ref{sec:featurecsv}) \\
Standard summary numbers (mean, variance, SE) & \texttt{-{}-classical} & Long CSV (\S\ref{sec:classicalcsv}) \\
Diversity along a chromosome & \texttt{-{}-classical -{}-rolling-window} & Long CSV, one block per window (\S\ref{sec:flags}) \\
Drop-in R / P / C files & \texttt{-{}-adze-prefix} & R/P/C text files (\S\ref{sec:rpc}) \\
Per-locus values & \texttt{-{}-adze-full} or \texttt{-{}-per-locus-out} & FULL / long files (\S\ref{sec:full}) \\
\bottomrule
\end{tabular}
\end{center}

\subsection{Feature CSV (default)}
\label{sec:featurecsv}

A wide CSV with \textbf{one row per rarefaction depth \(g\)}. For every statistic
there are \textbf{five moment columns}, named \texttt{<statistic>\_<moment>},
where \texttt{<moment>} is one of \texttt{mean}, \texttt{variance}, \texttt{se},
\texttt{skewness}, \texttt{kurtosis} (excess kurtosis). The statistics are:
\begin{itemize}
  \item \texttt{alpha\_j}---allelic richness for population \(j\).
  \item \texttt{pi\_j}---private allelic richness for population \(j\).
  \item \texttt{pihat\_XY}---combination-private richness for the population
        combination \(X,Y\) (for example \texttt{pihat\_12}, \texttt{pihat\_13},
        \texttt{pihat\_23}); combination size is set by \texttt{-{}-pihat-sizes}.
\end{itemize}

For the three-population example that is \texttt{alpha\_1..alpha\_3},
\texttt{pi\_1..pi\_3}, and \texttt{pihat\_12}, \texttt{pihat\_13},
\texttt{pihat\_23}---nine statistics \(\times\) five moments = 45 columns, plus
the leading \texttt{g} column, giving 46. \texttt{-{}-moments} restricts which of
the five suffixes are written.

\subsection{Classical CSV (\texttt{-{}-classical})}
\label{sec:classicalcsv}

A long-format table of the three classical moments, one row per
\texttt{(statistic, depth)}:
\begin{lstlisting}[style=out]
statistic,g,n_loci,mean,variance,se
\end{lstlisting}

\subsection{R / P / C Files (\texttt{-{}-adze-prefix})}
\label{sec:rpc}

\texttt{-{}-adze-prefix PREFIX} writes three summary files, one grouping per
block:
\begin{itemize}
  \item \texttt{PREFIX\_r}---allelic richness (R file).
  \item \texttt{PREFIX\_p}---private allelic richness (P file).
  \item \texttt{PREFIX\_c\_K}---combination-private richness for combination
        size \(K\) (C file).
\end{itemize}
Each data line is \texttt{grouping \ldots{} g num\_loci mean variance std\_err}.

\subsection{FULL Per-Locus Files (\texttt{-{}-adze-full})}
\label{sec:full}

Adding \texttt{-{}-adze-full} to \texttt{-{}-adze-prefix} also writes per-locus
files:
\begin{itemize}
  \item \texttt{PREFIX\_r\_fulldata}
  \item \texttt{PREFIX\_p\_fulldata}
  \item \texttt{PREFIX\_c\_K\_fulldata}
\end{itemize}
Each row carries the grouping, depth \(g\), locus count, \textbf{one column per
locus} (the per-locus statistic value), then the \texttt{AVG}, \texttt{VAR}, and
\texttt{STD\_ERR} summary columns. \texttt{-{}-per-locus-out PATH} writes the
same per-locus values in a single long-format CSV.

\subsection{Deleted-Loci Report}
\label{sec:deleted}

When \texttt{-{}-max-missing-per-population} drops loci in R/P/C mode, \padze{}
writes a companion report \texttt{PREFIX\_p\_deletedloci} listing the loci
excluded because at least one grouping exceeded the missing tolerance.

% ===========================================================================
\section{Worked Example}
\label{sec:example}

This end-to-end run turns the shipped VCF into a complete set of R/P/C output
files. From the \texttt{GitHub/} directory:
\begin{lstlisting}[style=cmd]
PYTHONPATH=src python -m padze features \
  --vcf examples/trio.vcf \
  --popmap examples/trio.popmap \
  --population-order A B C \
  --adze-prefix out \
  --adze-full
\end{lstlisting}

\padze{} reports the files it wrote:
\begin{lstlisting}[style=out]
wrote ADZE-compatible files: out_r, out_r_fulldata, out_p, out_p_fulldata, out_c_2, out_c_2_fulldata
\end{lstlisting}
\begin{itemize}
  \item \texttt{out\_r}, \texttt{out\_p}, \texttt{out\_c\_2}---the R, P, and C
        summary files.
  \item \texttt{out\_r\_fulldata}, \texttt{out\_p\_fulldata},
        \texttt{out\_c\_2\_fulldata}---the FULL per-locus files.
\end{itemize}

The R summary file \texttt{out\_r} (allelic richness per population and depth):
\begin{lstlisting}[style=out]
A 2 7 1.380952380952381 0.07142857142857144 0.10101525445522108
A 3 7 1.5714285714285714 0.1607142857142857 0.1515228816828316

B 2 7 1.3095238095238098 0.08730158730158731 0.11167656571008165
B 3 7 1.4642857142857142 0.19642857142857142 0.16751484856512247
B 4 7 1.5714285714285714 0.2857142857142857 0.20203050891044214

C 2 7 1.4761904761904763 0.050264550264550255 0.08473871628596277
C 3 7 1.7142857142857142 0.11309523809523814 0.1271080744289442
C 4 7 1.8571428571428572 0.1428571428571429 0.14285714285714288
\end{lstlisting}

Each block reports one population; columns are grouping, depth \(g\), number of
loci, mean, variance, and standard error. Each population is carried to the
largest depth its own sample size supports---\(B\) and \(C\) keep all four gene
copies at every locus and reach \(g = 4\), while \(A\) drops to \(g = 3\) because
of the missing call at \texttt{chr1:600} (Section~\ref{sec:depth}).

The FULL R file \texttt{out\_r\_fulldata} adds one column per locus. Its header
and the population-\(A\) block:
\begin{lstlisting}[style=out]
POP_GROUPING G NUM_LOCI chr1:100 chr1:200 chr1:300 chr1:400 chr1:500 chr1:600 chr1:700 AVG VAR STD_ERR
A 2 7 1.5 1.5 1.5 1.5 1 1.6666666666666667 1 1.380952380952381 0.07142857142857144 0.10101525445522108
A 3 7 1.75 1.75 1.75 1.75 1 2 1 1.5714285714285714 0.1607142857142857 0.1515228816828316
\end{lstlisting}

The C file \texttt{out\_c\_2} reports combination-private richness for each pair
of populations:
\begin{lstlisting}[style=out]
A B 2 7 0.32142857142857145 0.026455026455026457 0.061475926130276456
A B 3 7 0.24107142857142858 0.0630580357142857 0.09491201618513079

A C 2 7 0.4761904761904762 0.13128306878306878 0.13694788830743965
A C 3 7 0.39285714285714285 0.18452380952380953 0.16235930591649828

B C 2 7 0.2976190476190476 0.08531746031746033 0.11040022018447268
B C 3 7 0.24107142857142858 0.1255580357142857 0.13392857142857142
\end{lstlisting}

% ===========================================================================
\section{Validation and Correctness}
\label{sec:validation}

\padze{} is a validated implementation of the rarefaction method.
\begin{itemize}
  \item \textbf{Automated test suite.} 82 automated tests pass (with 11
        environment-gated tests skipped when optional external artifacts are
        absent). Run them with \texttt{pytest} after
        \texttt{pip install -e ".[test]"}.
  \item \textbf{Numerical core validated against the original C++ ADZE.} The
        rarefaction core, the classical mean / variance / standard-error
        summaries, and the per-locus FULL outputs are checked against the
        original C++ ADZE and agree to print precision.
  \item \textbf{VCF path validated against STRUCTURE.} The VCF reader is
        validated against matched STRUCTURE encodings of the same
        genotypes---including allele-coding permutations, multiallelic sites, and
        missing-data edge cases---so the VCF input path yields exactly the
        STRUCTURE-derived counts and statistics. A STRUCTURE encoding of the trio
        example reproduces byte-for-byte the classical values shown in
        Section~\ref{sec:qs-classical}.
\end{itemize}

For the full verification procedure---commands, expected shapes, and the
cross-tool comparison---see \texttt{docs/usage-and-verification.md}.

% ===========================================================================
\section{References}
\label{sec:references}

If you use \padze{}'s allelic-rarefaction statistics, cite the paper that
introduced the method:

\begin{quote}
Szpiech ZA, Jakobsson M, Rosenberg NA. 2008. ADZE: a rarefaction approach for
counting alleles private to combinations of populations. \textit{Bioinformatics}
24:2498--2504.
doi:\href{https://doi.org/10.1093/bioinformatics/btn478}{10.1093/bioinformatics/btn478}.
\end{quote}

Original ADZE 1.0 source, manual, and prebuilt binaries:
\url{https://github.com/szpiech/ADZE}.

\padze{} is an independent Python implementation, not endorsed by or affiliated
with the original ADZE authors, and is distributed under the MIT license.

\end{document}
